\documentclass[11pt]{article}

\usepackage[margin=1in]{geometry}
\usepackage{booktabs}
\usepackage{amsmath}
\usepackage{amssymb}
\usepackage{textcomp}
\usepackage{hyperref}
\usepackage[symbol]{footmisc}

% Convenience macros for the two methods, so wording stays consistent.
\newcommand{\dnlp}{the DNLP method}
\newcommand{\std}{the standard (Pypower/pips) method}

\title{A comparison of disciplined nonlinear programming and\\
Pypower on AC optimal power flow}
\author{Bennet Meyers}
\date{\today}

\begin{document}
\maketitle

\section{Summary}

We compared two methods for solving the same AC optimal power flow (AC-OPF)
problems on the standard 9-bus network:

\begin{enumerate}
  \item \textbf{The DNLP method} --- the problem is expressed in CVXPY and
    canonicalized through its disciplined nonlinear programming (DNLP)
    framework before being handed to the interior-point solver IPOPT.
  \item \textbf{The standard (Pypower/pips) method} --- the same physical
    problem is handed directly to Pypower's native primal-dual interior-point
    solver (pips).
\end{enumerate}

Both methods target the identical nonconvex AC-OPF (the same sinusoidal power
flow physics, the same generator and voltage limits). The only difference under
test is \emph{how the problem is presented to the interior-point method}: the
DNLP method canonicalizes it first, the standard method does not.

We varied two features of the problem, each on or off:

\begin{itemize}
  \item \textbf{Generator costs}. The base case prices generators with smooth,
    quadratic cost curves. The test case swapped the three generator curves for
    piecewise-linear functions.
  \item \textbf{HVDC transmission lines}. The base case had no HVDC line links.
    The test case added three controllable DC links, each moving real power
    between two buses subject to capacity limits and a proportional transport
    loss.
\end{itemize}

The headline result: \textbf{each feature alone is harmless --- both methods
agree to four decimal places --- but when both features are present together,
the standard method settles for a solution roughly 14\% more expensive than the
DNLP method.} The common thread is not nonconvexity (both added features are
convex); it is \textbf{nondifferentiability}. Piecewise-linear costs have kinks, and the optimal DC-line 
dispatch sits at capacity corners. An interior-point method
models the objective and constraints with local gradients and curvature, which
are ill-defined at a kink or corner. The DNLP method instead rewrites each
nondifferentiable convex piece into a smooth (indeed linear) epigraph form that
the interior-point method can descend exactly.

\section{Methods}

\subsection{Network models}

The base network is the textbook 9-bus, 3-generator case. On top of it we toggle
two independent features, giving four problem variants:

\begin{center}
\begin{tabular}{lll}
\toprule
 & \textbf{No DC lines} & \textbf{With DC lines} \\
\midrule
\textbf{Smooth (quadratic) costs} & plain 9-bus & 9-bus + 3 HVDC links \\
\textbf{Piecewise-linear costs} & 9-bus, PWL costs & 9-bus + PWL costs + 3 HVDC links \\
\bottomrule
\end{tabular}
\end{center}

\paragraph{Smooth vs.\ piecewise-linear costs.} In the smooth variant every
generator has a quadratic cost $c_2 P^2 + c_1 P + c_0$. In the piecewise-linear
variant two of the three generators are instead priced with a convex
piecewise-linear curve --- a sequence of straight-line segments of increasing
slope. A piecewise-linear convex cost is still convex, but it is
not differentiable at the breakpoints where two segments meet.

\paragraph{The three HVDC links.} When present, the three DC links (with their
from-bus $\to$ to-bus, capacity range, and proportional loss) are:

\begin{center}
\begin{tabular}{cccc}
\toprule
link & from $\to$ to & capacity (MW) & transport loss \\
\midrule
A & bus 30\footnotemark $\to$ bus 4 & $[1, 10]$ & 1\% (plus a 
	small fixed loss) \\
B & bus 7 $\to$ bus 9 & $[2, 10]$ & 0\% (lossless) \\
C & bus 5 $\to$ bus 9 & $[0, 10]$ & 5\% \\
\bottomrule
\end{tabular}
\end{center}
\footnotetext{Note that this test case has a ``random'' bus ID
	number (bus 30), which is either a transcription error from back in the day or
	an explicit test to make sure people track bus IDs correctly.}

A DC link is modelled as a controllable real-power injection: it withdraws power
at its from-bus and delivers it (minus the transport loss) at its to-bus. Two of
the three links (B and C) deliver into the \textbf{same} bus (bus 9) but
withdraw from \textbf{different} buses (7 and 5), which becomes important below.



\subsection{Model agreement between languages}

To compare the two approaches on genuinely the same problem, we neutralize a few
incidental modelling differences so that only the solution method differs:

\begin{itemize}
  \item \textbf{Branch flow limits are removed} on both sides (set to a very
    large value). cvxopf does not yet enforce branch limits, so we relax
    Pypower's as well.
  \item \textbf{DC-line reactive power is pinned to zero} on both sides (the DC
    links are modelled at unity power factor), and their terminal buses are
    treated identically.
  \item \textbf{The fixed (no-load) converter loss on link A is applied
    identically on both sides.} cvxopf normally omits this small fixed loss;
    here we add it back with a single linear constraint (valid because the flow
    direction on link A is known in advance), so neither approach is given an
    artificial advantage.
\end{itemize}

With these in place, the two solve paths address the identical optimization
problem, and any difference in the answer is attributable to the solution method
alone. We verified separately that the two sides build \textbf{the same
electrical network} (their bus admittance matrices agree to machine precision)
and that no DC-line endpoint was transposed between the two model
specifications.

\subsection{Analysis}

For each variant we compare the two methods on:

\begin{itemize}
  \item \textbf{Objective value} (total generation cost).
  \item \textbf{Generator dispatch} (how much each generator produces).
  \item \textbf{DC-line flows} (how power is routed through the links).
\end{itemize}

When the two disagree, we test \emph{why} by evaluating one method's solution
inside the other method's problem, without re-solving. If a solution is fully
feasible in the other's problem yet cheaper than what that method found, the
more expensive answer is demonstrably suboptimal.

\subsection{Reproducing the results}

The \texttt{demo.py} script alongside this report reproduces every number in it.
It builds the four problem variants, solves each with the DNLP method live,
compares against the committed Pypower reference, and prints the $2\times2$
table, the suboptimality proof, the routing decomposition, and the marginal-cost
exhibit.

From the repository root, in the main cvxopf environment:

\begin{verbatim}
uv run --active python experiments/dnlp_vs_pypower/demo.py
\end{verbatim}

The Pypower reference solutions are precomputed and committed
(\texttt{reference\_pypower.json}) because cvxopf and Pypower require separate
Python environments and cannot share a process; the two no-DC cells are copied
from the repository's test-validated fixtures and the two DC cells are computed
with the neutralized dcline model. To regenerate the reference in its isolated
Pypower sandbox (not needed to run the demo):

\begin{verbatim}
uv run experiments/dnlp_vs_pypower/generate_pypower_reference.py
\end{verbatim}

\section{Observations}

\paragraph{Plain 9-bus, smooth costs.} The two methods agree to within 0.01\% on
objective and to a fraction of a MW on every generator. This is the expected
baseline --- a smooth, well-behaved AC-OPF that a standard interior-point method
handles well. This is an official test case in the cvxopf repository, one of the
first tests written in fact.

\paragraph{9-bus with piecewise-linear costs, no DC lines.} The two methods
again agree to within 0.01\% on objective and on every generator's output.
Piecewise-linear costs alone do not force a divergence on this network. This is
also a test case in the cvxopf repository and was written when PWL costs were
added.

\paragraph{9-bus with DC lines, smooth costs.} Here the generator dispatch of
the two methods agrees closely --- total generator output differs by under 2 MW,
against a 315 MW load --- but the \textbf{DC-line routing differs sharply}.
cvxopf runs the lossless link B at its floor (2 MW) and the lossy link C at its
ceiling (10 MW); Pypower does the reverse, maxing the lossless link B (10 MW)
and idling the lossy link C ($\approx 0$ MW). cvxopf's objective is about 1\%
lower. This difference is investigated in more detail below, but the two methods
achieve a very similar objective cost and generator dispatch.

\paragraph{9-bus with both piecewise-linear costs and DC lines.} Now the two
methods diverge dramatically. cvxopf finds a solution costing about 5469;
Pypower settles at about 6250, \textbf{roughly 14\% more expensive}. The
generator dispatch differs by over 130 MW, and the DC-line routing is again
inverted between the two.

The standard method's answer is a genuine suboptimality, not a modelling
artifact. The DNLP method's cheaper solution is \textbf{fully feasible in the
standard formulation}, every power balance and operating limit is satisfied to
machine precision. In other words, cvxpy and IPOPT returned a feasible solution
to the Pypower problem that has a lower objective value. We can't prove the
cvxopf solution is the globally optimal point, but we've proved that the Pypower
solution is not.

\subsection{Discrepancies on the ``DC line, smooth cost'' case}

The routing flip looks backwards at first. Why would the DNLP method pay a 5\%
transport loss on link C instead of using the free link B? The ${\sim}1\%$ cost
gap is the sum of two separable effects, and pinning both methods to the
\emph{same} routing lets us isolate each.

\paragraph{A small, genuine cost optimum ($\approx 18$ units).} Links B and C
deliver to the same bus but withdraw from different buses, so the routing choice
changes \emph{where} on the network power is drawn, and therefore which
generators cover it. This is a generation-cost effect, not a loss effect, which
we establish in two steps. First, we rule out losses directly. The DNLP routing
runs the 5\% lossy link on purpose, so its DC transport loss is higher; this
outweighs its slightly lower AC network loss, leaving its \textbf{total} loss
\emph{higher} than the standard routing's (4.87 vs 4.64 MW). If the DNLP routing
had won by saving losses it would be the more expensive solution, not the
cheaper one. Second, we confirm the advantage is in the generation mix by
comparing marginal costs --- the slope of each generator's cost curve at its
operating point, i.e.\ the cost of its next MW. At the DNLP dispatch these
slopes sit closer together across the generators not pinned at a limit; at the
standard routing they are more spread out. A tighter spread means output has
shifted toward the cheaper margin, which is why the DNLP routing costs about 18
units less (on about 0.77 MW less total generation). This is a real but small
optimum on a nearly flat part of the landscape. (The marginal-cost comparison is
corroborating evidence, not a proof of optimality: the full AC-OPF is nonconvex,
so there is no clean theorem that equal marginal costs mark the optimum. The
loss decomposition is what rules out the competing loss-based explanation.)

\paragraph{The standard method is stuck ($\approx 34$ units).} Independently of
which routing is better, at the standard method's \emph{own} routing cvxopf
reaches a solution about 34 units cheaper than Pypower. The standard method is
suboptimal even at the operating point it itself selected --- it has not fully
descended its own problem.

Both effects favour the DNLP method, and together they account for the
${\sim}1\%$ gap in this cell even though the generator dispatch otherwise matches
closely.

\section{Analysis}

\subsection{The pattern}

\begin{center}
\begin{tabular}{lll}
\toprule
variant & added nondifferentiable structure & result \\
\midrule
plain 9-bus, smooth & none & exact agreement \\
PWL costs only & cost kinks & exact agreement \\
DC lines only & capacity-corner routing & dispatch agrees; ${\sim}1\%$ routing gap, DNLP cheaper \\
PWL costs \textbf{and} DC lines & kinks \textbf{and} corners & ${\sim}14\%$ gap, DNLP cheaper \\
\bottomrule
\end{tabular}
\end{center}

Neither added feature causes (much) trouble on its own. Together they do. The
gap grows with the amount of nondifferentiable convex structure in the problem,
and every time the two methods disagree, the DNLP method is the one that finds
the cheaper solution. We have verified in each case that the standard method's
answer is genuinely suboptimal --- a feasible, cheaper point exists in its own
feasible set.

\subsection{Nondifferentiability is the core issue}

Both added features are \textbf{convex}; what they introduce is
\textbf{nondifferentiability}. Piecewise-linear costs have kinks at their
breakpoints, and the optimal DC-line dispatch sits at a capacity corner. The
underlying AC power flow is nonconvex, but that nonconvexity is common to both
methods and is not where they diverge --- they agree perfectly on the smooth
variants, which share the same AC nonconvexity.

An interior-point method builds a local model from first and second derivatives.
At a kink or an active-set corner those derivatives are not well-defined, so the
local model is wrong precisely where the optimum lives and the method can stall
short, which is what we appear to see. The DNLP method avoids this by rewriting
each nondifferentiable convex atom into its \textbf{epigraph form} before the
solver sees it: a piecewise-linear cost becomes the maximum of its affine
segments (one auxiliary variable plus linear inequalities), giving the solver a
smooth landscape where its gradient-and-curvature model is valid everywhere.

We present this as the \emph{likely} mechanism rather than a measured one. It is
grounded in the construction (the piecewise-linear cost is built directly as a
pointwise maximum of affine segments), but we did not instrument the canonical
program at runtime to confirm it. The empirical claims (exact agreement on the
single-feature variants, ${\sim}14\%$ divergence combined, verified
suboptimality) do not depend on this account.

\section{Conclusion}

Across the four variants, the two methods returned identical solutions on the
plain 9-bus case and on each single-feature variant (piecewise-linear costs
alone, and DC lines alone with matching generator dispatch), and returned
different solutions only when piecewise-linear costs and DC lines were present
together, where the standard method's objective was about 14\% higher. In every
case where the two disagreed, the DNLP solution was the cheaper one, and we
verified that it was fully feasible in the standard method's own problem --- so
the standard method's more expensive answer is a suboptimal point of its own
feasible set, not the result of a different model. For the smooth-cost DC-line
case we further decomposed the ${\sim}1\%$ gap into a ${\sim}18$-unit
generation-cost difference and a ${\sim}34$-unit residual by which the standard
method is suboptimal even at its own routing.

What these results do not establish is a general guarantee: both problems are
nonconvex, we did not certify global optimality of any DNLP solution, and the
comparison covers one network at four operating points. What they do establish
is narrow and firm --- on these problems the two methods diverge only when
convex nondifferentiable structure is present, and where they diverge the
standard method returns a verifiably suboptimal solution.

\end{document}
